\input{preamble/preamble.tex}

\geometry{margin=0.85in}
\AtBeginDocument{\small}

\newcommand{\ExamNameField}{\noindent\textbf{Name:}\ \rule{0.7\linewidth}{0.4pt}\par\medskip}

\begin{document}
	\ExamTitleBlock{11th grade}{Learning evidence 2.1: Confidence Interval Planning (Solutions)}{}

	\ExamSection{Problems}
	\begin{ExamProblems}
		\item
		\subsection*{Problem description}
		A city budgeting office wants to estimate the average weekly transit subsidy per rider for two pilot zones.
		Zone A recorded the following sample of weekly subsidy amounts (USD):

		\begin{itemize}
			\item Zone A sample (\(n_A = 6\)): 18, 22, 19, 21, 20, 17.
			\item Zone B sample (\(n_B = 8\)): 24, 26, 23, 25, 27, 22, 24, 26.
		\end{itemize}

		An annual audit provides the population standard deviations for the zones: \(\sigma_A = 3.5\) and \(\sigma_B = 4.0\).
		For classroom purposes, suppose the true population mean weekly subsidy is \(\mu_A = 20\) dollars in Zone A and \(\mu_B = 25\) dollars in Zone B.
		Construct and interpret a 95\% confidence interval for the population mean weekly subsidy in each zone.

		\subsection*{C1}
		First, we calculate the sample means.
		\[
		\bar{x}_A = \frac{\sum x_A}{n_A}
		\]
		\[
		\sum x_A = 18 + 22 + 19 + 21 + 20 + 17 = 117
		\]
		\[
		\bar{x}_A = \frac{117}{6} = 19.5
		\]
		\[
		\bar{x}_B = \frac{\sum x_B}{n_B}
		\]
		\[
		\sum x_B = 24 + 26 + 23 + 25 + 27 + 22 + 24 + 26 = 197
		\]
		\[
		\bar{x}_B = \frac{197}{8} = 24.625
		\]
		Second, we calculate the sample standard deviations.
		\[
		s_A^2 = \frac{\sum (x_A-\bar{x}_A)^2}{n_A-1}
		\]
		\[
		\sum (x_A-\bar{x}_A)^2 = (18-19.5)^2 + (22-19.5)^2 + (19-19.5)^2 + (21-19.5)^2 + (20-19.5)^2 + (17-19.5)^2 = 17.5
		\]
		\[
		s_A = \sqrt{\frac{17.5}{6-1}} = \sqrt{3.5} \approx 1.87
		\]
		\[
		s_B^2 = \frac{\sum (x_B-\bar{x}_B)^2}{n_B-1}
		\]
		\[
		\sum (x_B-\bar{x}_B)^2 = (24-24.625)^2 + (26-24.625)^2 + (23-24.625)^2 + (25-24.625)^2 + (27-24.625)^2 + (22-24.625)^2 + (24-24.625)^2 + (26-24.625)^2 = 19.875
		\]
		\[
		s_B = \sqrt{\frac{19.875}{8-1}} \approx 1.69
		\]

		\subsection*{C2}
		The sample statistics are $\bar{x}_A = 19.5$, $s_A \approx 1.87$, $\bar{x}_B = 24.625$, and $s_B \approx 1.69$, while the population parameters are $\mu_A$, $\mu_B$, $\sigma_A = 3.5$, and $\sigma_B = 4.0$. The sample values estimate the population values.

		\subsection*{C3}
		A single point estimate can miss the true mean because of sampling variation. A confidence interval shows a range of plausible values for each $\mu$, which is more informative for planning.

		\subsection*{C5}
		With known $\sigma_A = 3.5$ and $n_A = 6$,
		\[
		SE_A = \frac{\sigma_A}{\sqrt{n_A}}
		\]
		\[
		SE_A = \frac{3.5}{\sqrt{6}} \approx 1.43
		\]
		For 95\% confidence, $z_{0.025}=1.96$.
		\[
		E_A = z_{0.025} \cdot SE_A
		\]
		\[
		E_A = 1.96(1.43) \approx 2.80
		\]
		\[
		\mu_A \in \bar{x}_A \pm E_A
		\]
		\[
		\mu_A \in 19.5 \pm 2.80 = [19.5 - 2.80, 19.5 + 2.80] = [16.70, 22.30]
		\]
		With known $\sigma_B = 4.0$ and $n_B = 8$,
		\[
		SE_B = \frac{\sigma_B}{\sqrt{n_B}}
		\]
		\[
		SE_B = \frac{4.0}{\sqrt{8}} \approx 1.41
		\]
		\[
		E_B = z_{0.025} \cdot SE_B
		\]
		\[
		E_B = 1.96(1.41) \approx 2.76
		\]
		\[
		\mu_B \in \bar{x}_B \pm E_B
		\]
		\[
		\mu_B \in 24.625 \pm 2.76 = [24.625 - 2.76, 24.625 + 2.76] = [21.87, 27.39]
		\]

		\subsection*{C4}
		We are 95\% confident the true mean weekly subsidy in Zone A is between about 16.70 and 22.30 dollars, and 95\% confident the true mean weekly subsidy in Zone B is between about 21.87 and 27.39 dollars. The intervals reflect uncertainty from sampling in each zone.

		\item
		\subsection*{Problem description}
		A regional warehouse studies the number of orders processed per hour to plan staffing.
		A random sample of 60 hours is summarized in grouped form below. The operations team also reports a known population standard deviation of \(\sigma = 6.2\) orders per hour.
		For academic purposes only, assume the true population mean order rate is \(\mu = 22\) orders per hour.
		The grouped data pairs (orders per hour, frequency) are: \(14 \rightarrow 8\), \(18 \rightarrow 16\), \(22 \rightarrow 20\), \(26 \rightarrow 10\), and \(30 \rightarrow 6\).
		Construct and interpret a 95\% confidence interval for the population mean orders per hour.

		\subsection*{C1}
		Using grouped values $x$ with frequencies $f$.
		\[
		\bar{x} = \frac{\sum f x}{n}
		\]
		\[
		n = \sum f
		\]
		\[
		n = 8 + 16 + 20 + 10 + 6 = 60
		\]
		\[
		\sum f x = 8(14) + 16(18) + 20(22) + 10(26) + 6(30) = 1280
		\]
		\[
		\bar{x} = \frac{1280}{60} \approx 21.33
		\]
		Second, we calculate the sample standard deviation $s$.
		\[
		s^2 = \frac{\sum f(x-\bar{x})^2}{n-1}
		\]
		\[
		\sum f(x-\bar{x})^2 \approx 8(14-21.33)^2 + 16(18-21.33)^2 + 20(22-21.33)^2 + 10(26-21.33)^2 + 6(30-21.33)^2
		\]
		\[
		\sum f(x-\bar{x})^2 \approx 430.22 + 177.78 + 8.89 + 217.78 + 450.67 = 1285.34
		\]
		\[
		s = \sqrt{\frac{\sum f(x-\bar{x})^2}{n-1}}
		\]
		\[
		s = \sqrt{\frac{1285.34}{60-1}} \approx 4.67
		\]

		\subsection*{C2}
		The sample statistics are $\bar{x} \approx 21.33$ and $s \approx 4.67$, while the population parameters are $\mu$ and $\sigma = 6.2$. The sample values estimate the population values.

		\subsection*{C3}
		The sample summarizes only 60 hours, so the true mean could differ from the point estimate. A confidence interval expresses the likely range for $\mu$ based on the data.

		\subsection*{C5}
		With known $\sigma = 6.2$ and $n=60$,
		\[
		SE = \frac{\sigma}{\sqrt{n}}
		\]
		\[
		SE = \frac{6.2}{\sqrt{60}} \approx 0.80
		\]
		For 95\% confidence, $z_{0.025}=1.96$.
		\[
		E = z_{0.025} \cdot SE
		\]
		\[
		E = 1.96(0.80) \approx 1.57
		\]
		\[
		\mu \in \bar{x} \pm E
		\]
		\[
		\mu \in 21.33 \pm 1.57 = [21.33 - 1.57, 21.33 + 1.57] = [19.76, 22.90]
		\]

		\subsection*{C4}
		We are 95\% confident the true mean orders processed per hour is between about 19.76 and 22.90 orders. The interval shows the range supported by the sample and known standard deviation.

		\item
		\subsection*{Problem description}
		A logistics firm monitors the time (in minutes) required to load containers at a port.
		A sample of 80 loading times is grouped into class intervals below. The firm has a historical estimate of the population standard deviation, \(\sigma = 5.8\) minutes.
		For classroom reference, suppose the true population mean loading time is \(\mu = 32\) minutes.
		The grouped intervals and frequencies are: 20--24 minutes (12), 25--29 minutes (18), 30--34 minutes (20), 35--39 minutes (17), and 40--44 minutes (13).
		Construct and interpret a 95\% confidence interval for the population mean loading time.

		\subsection*{C1}
		Using grouped midpoints $x$ with frequencies $f$.
		\[
		\bar{x} = \frac{\sum f x}{n}
		\]
		\[
		n = \sum f
		\]
		\[
		n = 12 + 18 + 20 + 17 + 13 = 80
		\]
		\[
		\sum f x = 12(22) + 18(27) + 20(32) + 17(37) + 13(42) = 2565
		\]
		\[
		\bar{x} = \frac{2565}{80} = 32.0625
		\]
		Second, we calculate the sample standard deviation $s$.
		\[
		s^2 = \frac{\sum f(x-\bar{x})^2}{n-1}
		\]
		\[
		\sum f(x-\bar{x})^2 \approx 12(22-32.0625)^2 + 18(27-32.0625)^2 + 20(32-32.0625)^2 + 17(37-32.0625)^2 + 13(42-32.0625)^2
		\]
		\[
		\sum f(x-\bar{x})^2 \approx 3374.69
		\]
		\[
		s = \sqrt{\frac{\sum f(x-\bar{x})^2}{n-1}}
		\]
		\[
		s = \sqrt{\frac{3374.69}{80-1}} \approx 6.54
		\]

		\subsection*{C2}
		The sample statistics are $\bar{x} = 32.0625$ and $s \approx 6.54$, while the population parameters are $\mu$ and $\sigma = 5.8$. The sample values estimate the population values.

		\subsection*{C3}
		Even with 80 observations, a point estimate does not show how much the mean could vary. A confidence interval summarizes that uncertainty by giving a plausible range for $\mu$.

		\subsection*{C5}
		With known $\sigma = 5.8$ and $n=80$,
		\[
		SE = \frac{\sigma}{\sqrt{n}}
		\]
		\[
		SE = \frac{5.8}{\sqrt{80}} \approx 0.65
		\]
		For 95\% confidence, $z_{0.025}=1.96$.
		\[
		E = z_{0.025} \cdot SE
		\]
		\[
		E = 1.96(0.65) \approx 1.27
		\]
		\[
		\mu \in \bar{x} \pm E
		\]
		\[
		\mu \in 32.0625 \pm 1.27 = [32.0625 - 1.27, 32.0625 + 1.27] = [30.79, 33.33]
		\]

		\subsection*{C4}
		We are 95\% confident the true mean loading time is between about 30.79 and 33.33 minutes. The interval provides a plausible range for the population mean based on the sample.
	\end{ExamProblems}
\end{document}
